(1) 设 \(a=Z\alpha\)、\(\kappa=j+\tfrac12\)，
\[
\sqrt{\kappa^2-a^2}=\kappa-\frac{a^2}{2\kappa}-\frac{a^4}{8\kappa^3}+O(a^6).
\]
分母 \(d=n-a^2/(2\kappa)+O(a^4)\)，故 \(a^2/d^2=a^2/n^2+a^4/(\kappa n^3)+O(a^6)\)。再用 \((1+x)^{-1/2}=1-x/2+3x^2/8+\cdots\)，得
\[
E_{nj}=mc^2-\frac{mc^2a^2}{2n^2}-\frac{mc^2a^4}{2n^4}\left(\frac n\kappa-\frac34\right)+O(a^6).
\]
第二项为Bohr能级，四阶项为相对论精细结构。
(2) \(n=2\) 时 \(j=1/2\) 的四阶项为 \(-5mc^2a^4/128\)，\(j=3/2\) 的四阶项为 \(-mc^2a^4/128\)。因此
\[
E_{2,3/2}-E_{2,1/2}=\frac{mc^2a^4}{32}+O(a^6)>0.
\]
能量只依赖 \(n,j\)，每组的 \(m_j=-j,\ldots,j\) 在无外场下仍简并。
(3) 主束缚能量级 \(|E-mc^2|\sim mc^2a^2/(2n^2)\)，精细结构相对比约为 \(a^2/n^2\) 乘角动量因子。\(Z\to0\) 时束缚与劈裂均消失；固定 \(n\) 时精细结构 \(\propto Z^4\)（在 \(Z\alpha\ll1\) 区间）。纯Coulomb Dirac哈密顿量具有额外对称性，能谱只依赖 \(n,j\) 而不独立依赖 \(l\)，故例如 \(2s_{1/2}\) 与 \(2p_{1/2}\) 在此理想模型中简并。
(4) 根号 \(\sqrt{(j+1/2)^2-a^2}\) 要求 \(a<j+1/2\)，最严的是 \(j=1/2\)，即 \(Z\alpha<1\)；等号是点核模型的临界边界。实际核有有限尺寸，且辐射修正如真空极化、自能Lamb位移会破除理想简并；高 \(Z\) 还需更完整强场QED处理，不能把上述小 \(a\) 四阶展开当作高核电荷精确答案。
